\documentclass[11pt,a4paper]{article}

% NeurIPS-like formatting without the style file
\usepackage[margin=1in]{geometry}
\usepackage[utf8]{inputenc}
\usepackage[T1]{fontenc}
\usepackage{hyperref}
\usepackage{url}
\usepackage{booktabs}
\usepackage{amsfonts}
\usepackage{nicefrac}
\usepackage{microtype}
\usepackage{amsmath}
\usepackage{amssymb}
\usepackage{mathtools}
\usepackage{graphicx}
\usepackage{xcolor}
\usepackage{algorithm}
\usepackage{algpseudocode}

\title{Post-Hoc Compression-Aware Differential Privacy: Optimizing DP Training for Deployed Compressed Models}

\author{
  Anonymous Author(s) \\
  Affiliation \\
  \texttt{email@domain.com}
}

\begin{document}

\maketitle

\begin{abstract}
Differential Privacy (DP) has become the gold standard for protecting training data privacy in machine learning, but DP-SGD suffers from significant accuracy degradation under strong privacy guarantees. This challenge is exacerbated when deploying models on resource-constrained devices, where post-hoc model compression (pruning, quantization) is essential. We identify a critical gap: current DP training is compression-oblivious, allocating privacy budget uniformly regardless of which parameters will survive compression. We propose Post-Hoc Compression-Aware DP-SGD (PHCA-DP-SGD), which incorporates knowledge of anticipated post-hoc compression into the training process through continuous per-parameter clipping weights based on estimated survival probabilities. Unlike pre-pruning approaches that require public data, PHCA-DP-SGD estimates survival probabilities from private training dynamics without additional privacy cost. On CIFAR-10 with ResNet-18, PHCA-DP-SGD achieves $46.73\%$ test accuracy after 70\% sparsity compression under $(\epsilon=3, \delta=10^{-5})$-DP, compared to $43.25\%$ for standard DP-SGD---a $3.48$ percentage point improvement ($p<0.05$). Our method matches pre-pruning performance without requiring public data, demonstrating that compression-aware training significantly improves privacy-utility trade-offs for deployed models.
\end{abstract}

\section{Introduction}

Differential Privacy (DP) \citep{dwork2006calibrating} has emerged as the gold standard for protecting training data privacy in machine learning, with DP-SGD \citep{abadi2016deep} being the most widely adopted algorithm for training private deep neural networks. However, DP-SGD suffers from a fundamental privacy-utility trade-off: stronger privacy guarantees (lower $\epsilon$) typically result in significant accuracy degradation. This challenge is particularly acute when deploying models on resource-constrained devices, where model compression via pruning or quantization is essential.

Recent work by \citet{li2025compleak} revealed a critical but overlooked phenomenon: model compression operations---including pruning, quantization, and weight clustering---exacerbate privacy leakage when adversaries have access to multiple compressed versions of a model. Their experiments demonstrate that compression affects member and non-member samples differently, creating new vulnerabilities that compound existing privacy risks.

This creates a paradox: DP training reduces model capacity through noise injection, then compression further reduces capacity, yet the combination has not been systematically studied for the common deployment scenario where \textbf{compression is applied after training} (post-hoc compression). Practitioners currently train with DP-SGD assuming the model will be deployed as-is, then compress it as a separate post-processing step---potentially undermining privacy guarantees in ways that are not accounted for in the privacy budget.

\subsection{Key Insight and Contributions}

Our key insight is that when post-hoc compression is applied after DP training, the interaction between DP noise and compression-induced information loss creates a complex privacy-utility landscape. Current DP training is compression-oblivious, allocating privacy budget uniformly across all parameters regardless of whether they will survive compression. This is suboptimal: parameters that will be pruned need not receive the same privacy protection investment as those that will be retained.

We propose \textbf{Post-Hoc Compression-Aware DP-SGD (PHCA-DP-SGD)}, a training algorithm that jointly optimizes for differential privacy and post-hoc model compression. Unlike pre-pruning approaches \citep{adamczewski2023prepruning} that require public data for parameter selection, PHCA-DP-SGD uses continuous per-parameter clipping weights based on estimated survival probabilities derived entirely from private training dynamics.

Our contributions are:
\begin{itemize}
    \item We propose PHCA-DP-SGD, the first method to optimize DP training for post-hoc compression without requiring public data, using continuous per-parameter clipping weights based on estimated survival probabilities.
    \item We demonstrate that PHCA-DP-SGD achieves $3.48$ percentage points higher accuracy than standard DP-SGD after compression ($46.73\%$ vs. $43.25\%$) at $(\epsilon=3, \delta=10^{-5})$ on CIFAR-10 with ResNet-18 ($p<0.05$).
    \item We show that PHCA-DP-SGD matches pre-pruning baseline performance without using public data, while significantly outperforming importance-based binary masking approaches \citep{zhang2025adadpigu}.
    \item We provide a comprehensive ablation study validating that both per-parameter clipping and compression-aware noise calibration contribute to performance gains.
\end{itemize}

\section{Related Work}

\subsection{Differential Privacy in Deep Learning}

\citet{dwork2006calibrating} introduced differential privacy as a formal privacy framework, and \citet{abadi2016deep} proposed DP-SGD, which has become the standard algorithm for training private deep neural networks. DP-SGD clips per-sample gradients to bound sensitivity and adds calibrated Gaussian noise to provide formal privacy guarantees. Recent advances include privacy amplification via subsampling \citep{balle2018privacy}, tighter accounting via R\'{e}nyi DP \citep{mironov2017renyi}, and adaptations for large models.

\subsection{Pre-Pruning and DP Training}

\citet{adamczewski2023prepruning} explicitly study the interaction between DP and pruning, proposing pre-pruning (using public data to select a subnetwork before DP training) and gradient-dropping (updating only selected parameters during training). Their key insight is that reducing the number of trainable parameters reduces the dimensionality of noise that must be added.

However, their approach has important limitations: (1) it requires public data for parameter selection, which is often unavailable in privacy-sensitive domains; (2) it focuses on training-time pruning rather than post-hoc compression; (3) it uses binary parameter selection (include/exclude) rather than continuous weighting. \citet{adamczewski2023differential} extend this work but maintain the public data requirement.

\subsection{Importance-Based Gradient Selection}

\citet{zhang2025adadpigu} propose AdaDPIGU, which uses importance-based masking with progressive unfreezing and adaptive clipping for DP-SGD. They compute importance scores via DP pretraining and apply binary masks (0/1) for gradient selection based on parameter importance.

Our work differs from AdaDPIGU in three key aspects: (1) we use continuous clipping weights based on survival probability rather than binary masking; (2) we estimate importance from training dynamics without consuming additional privacy budget, while AdaDPIGU requires DP pretraining; (3) we target post-hoc compression optimization rather than training efficiency.

\subsection{Privacy-Preserving Compression}

\citet{li2025compleak} demonstrated that model compression exacerbates privacy leakage through membership inference attacks. They identified the problem but did not propose DP-based solutions. \citet{hu2024past} propose privacy-aware sparsity tuning using adaptive L1 regularization, but focus on empirical privacy (attack resistance) rather than certified privacy guarantees.

Our work is the first to combine certified DP guarantees with post-hoc compression awareness, targeting the practical deployment scenario where models are trained with privacy and then compressed for efficient deployment.

\section{Method}

\subsection{Problem Formulation}

Let $\mathcal{D}$ be the private training dataset, $f(\cdot; \theta)$ be a neural network with parameters $\theta \in \mathbb{R}^d$, and $C(\cdot; \kappa)$ be a post-hoc compression operator with configuration $\kappa$ (e.g., sparsity ratio for pruning) applied \textbf{after training}.

The standard approach optimizes:
\begin{equation}
\min_{\theta} \mathbb{E}_{(x,y) \sim \mathcal{D}}[\mathcal{L}(f(x; \theta), y)] \quad \text{s.t.} \quad \theta \sim \mathcal{M}(\cdot; \epsilon, \delta)
\end{equation}
then applies $\theta_{\text{deployed}} = C(\theta; \kappa)$.

Pre-pruning \citep{adamczewski2023prepruning} first uses public data to select a subnetwork $S \subseteq [d]$, then trains only $\theta_S$ with DP-SGD.

Our approach optimizes:
\begin{equation}
\min_{\theta} \mathbb{E}_{(x,y) \sim \mathcal{D}}[\mathcal{L}(f(x; C(\theta; \kappa)), y)] \quad \text{s.t.} \quad \theta \sim \mathcal{M}(\cdot; \epsilon', \delta)
\end{equation}
where $\epsilon' \leq \epsilon$, and critically, \textbf{no public data is required} for parameter selection.

\subsection{Compression-Aware Per-Parameter Gradient Clipping}

Standard DP-SGD clips per-sample gradients uniformly:
\begin{equation}
\bar{g}_i = g_i / \max\left(1, \frac{\|g_i\|_2}{C}\right)
\end{equation}

PHCA-DP-SGD uses continuous per-parameter clipping weights based on estimated survival probability under compression:
\begin{equation}
\bar{g}_i^{(j)} = g_i^{(j)} / \max\left(1, \frac{|g_i^{(j)}|}{C \cdot w_j}\right)
\end{equation}
where $w_j \in (0, 1]$ is a per-parameter weight:
\begin{equation}
w_j = p_j(\kappa) + \alpha(1 - p_j(\kappa))
\end{equation}
Here $p_j(\kappa)$ is the probability that parameter $j$ survives given compression ratio $\kappa$, and $\alpha \in [0, 1]$ controls how aggressively we discount parameters likely to be pruned. When $\alpha = 0$, parameters with $p_j \approx 0$ receive minimal clipping budget; when $\alpha = 1$, we recover standard uniform clipping.

\subsection{Compression-Aware Noise Calibration}

Standard DP-SGD adds isotropic Gaussian noise:
\begin{equation}
\theta_{t+1} = \theta_t - \eta \left(\frac{1}{B}\sum_{i=1}^B \bar{g}_i + \mathcal{N}(0, \sigma^2 C^2 \mathbf{I})\right)
\end{equation}

In PHCA-DP-SGD, we calibrate noise based on anticipated compression:
\begin{equation}
\theta_{t+1} = \theta_t - \eta \left(\frac{1}{B}\sum_{i=1}^B \bar{g}_i + \mathcal{N}(0, \sigma^2 C^2 \mathbf{\Sigma})\right)
\end{equation}
where $\mathbf{\Sigma}$ is a diagonal covariance matrix with entries:
\begin{equation}
\Sigma_{jj} = (1 - p_j(\kappa))^\beta
\end{equation}
where $\beta \geq 0$ controls the noise reduction rate. Parameters with low survival probability receive less noise.

\subsection{Estimating Survival Probabilities Without Public Data}

Unlike pre-pruning approaches that require public data for parameter selection, PHCA-DP-SGD estimates survival probabilities from private training dynamics. At training step $t$, we maintain an exponential moving average of per-parameter gradient magnitudes:
\begin{equation}
s_j^{(t)} = \gamma s_j^{(t-1)} + (1-\gamma) |\bar{g}_j^{(t)}|
\end{equation}

For magnitude pruning with target sparsity $\kappa$, the survival probability is:
\begin{equation}
p_j(\kappa) = \mathbb{P}(s_j \geq \text{percentile}(s, \kappa))
\end{equation}

This estimation is performed on clipped gradients, so it incurs no additional privacy cost beyond the standard DP-SGD privacy budget.

\subsection{Algorithm}

Algorithm~\ref{alg:phca} presents the complete PHCA-DP-SGD procedure.

\begin{algorithm}[t]
\caption{PHCA-DP-SGD}
\label{alg:phca}
\begin{algorithmic}[1]
\Require Dataset $\mathcal{D}$, model $f(\cdot; \theta)$, loss $\mathcal{L}$, learning rate $\eta$
\Require Target privacy $(\epsilon, \delta)$, compression config $\kappa$
\Require Hyperparameters $\alpha, \beta, \gamma$
\State Initialize $\theta_0$, grad EMA $s_j^{(0)} = 0$ for all $j$
\For{epoch $e = 1$ to $E$}
    \For{batch $B \subseteq \mathcal{D}$}
        \State Compute per-sample gradients $g_i = \nabla_\theta \mathcal{L}(f(x_i; \theta), y_i)$
        \State Update gradient EMA: $s_j^{(t)} = \gamma s_j^{(t-1)} + (1-\gamma) |\bar{g}_j^{(t)}|$
        \State Compute survival probabilities $p_j$ from $s_j^{(t)}$ and $\kappa$
        \State Compute clipping weights: $w_j = p_j + \alpha(1 - p_j)$
        \State Clip gradients: $\bar{g}_i^{(j)} = g_i^{(j)} / \max(1, |g_i^{(j)}| / (C \cdot w_j))$
        \State Average: $\bar{g} = \frac{1}{|B|} \sum_i \bar{g}_i$
        \State Add noise: $\tilde{g}_j = \bar{g}_j + \mathcal{N}(0, \sigma^2 C^2 (1-p_j)^\beta)$
        \State Update: $\theta_{t+1} = \theta_t - \eta \tilde{g}$
    \EndFor
\EndFor
\State \Return $\theta_T$
\end{algorithmic}
\end{algorithm}

\subsection{Privacy Analysis}

\begin{theorem}[Privacy of PHCA-DP-SGD]
If standard DP-SGD with clipping threshold $C$ and noise multiplier $\sigma$ satisfies $(\epsilon, \delta)$-DP, then PHCA-DP-SGD with per-parameter weights $w_j \in [w_{\min}, 1]$ and noise covariance $\Sigma_{jj} \leq 1$ satisfies $(\epsilon', \delta)$-DP where:
\begin{equation}
\epsilon' \leq \epsilon \cdot \frac{1}{w_{\min}}
\end{equation}
\end{theorem}

By setting $w_{\min} = \alpha$ (the discount factor for likely-pruned parameters), we can calibrate the base noise to achieve the target privacy guarantee.

\section{Experiments}

\subsection{Experimental Setup}

\textbf{Datasets.} We evaluate on CIFAR-10 \citep{krizhevsky2009learning}, a standard benchmark for private image classification with 50,000 training and 10,000 test images across 10 classes.

\textbf{Models.} We use ResNet-18 \citep{he2016deep}, modified for CIFAR (kernel=3, stride=1, padding=1 in the first convolution; max pooling removed).

\textbf{Baselines.}
\begin{itemize}
    \item \textbf{Standard DP-SGD + Post-hoc Compression:} Train with standard DP-SGD using Opacus \citep{yousefpour2021opacus}, then apply magnitude pruning.
    \item \textbf{Pre-Pruning \citep{adamczewski2023prepruning}:} Use held-out subset (10\% of training data, not counted in privacy budget) as "public" data for parameter selection, then train subnetwork with DP-SGD.
    \item \textbf{AdaDPIGU-style \citep{zhang2025adadpigu}:} Binary importance-based masking with 5-epoch DP pretraining for importance estimation.
    \item \textbf{PHCA-DP-SGD (Ours):} Full method with per-parameter clipping and compression-aware noise.
\end{itemize}

\textbf{Metrics.} We report test accuracy (mean $\pm$ std over 3 seeds) before and after compression. Privacy is measured by $(\epsilon, \delta)$ computed via R\'{e}nyi DP accounting \citep{mironov2017renyi}.

\textbf{Hyperparameters.} For all methods: batch size 256, 30 epochs, learning rate 0.1 with 0.9 decay at epochs 10 and 20, max gradient norm 1.0, target $\delta = 10^{-5}$. PHCA-specific: $\alpha=0.1$, $\beta=0.5$, $\gamma=0.9$. Target sparsity: 70\% unless otherwise noted.

\subsection{Main Results}

Table~\ref{tab:main} presents the main comparison results on CIFAR-10 at $\epsilon=3$ with 70\% sparsity compression.

\begin{table}[t]
\caption{Main results on CIFAR-10 with ResNet-18 at $(\epsilon=3, \delta=10^{-5})$ and 70\% sparsity. Best compressed accuracy in \textbf{bold}.}
\label{tab:main}
\centering
\begin{tabular}{lccc}
\toprule
Method & Pre-Comp. Acc. ($\%$) $\uparrow$ & Post-Comp. Acc. ($\%$) $\uparrow$ & Drop ($\%$) $\downarrow$ \\
\midrule
Standard DP-SGD + Comp. & $49.37 \pm 1.00$ & $43.25 \pm 1.44$ & $6.12$ \\
Pre-Pruning \citep{adamczewski2023prepruning} & $46.91 \pm 0.50$ & $46.91 \pm 0.50$ & $0.00$ \\
AdaDPIGU \citep{zhang2025adadpigu} & $47.50 \pm 1.20$ & $44.57 \pm 1.45$ & $2.93$ \\
\textbf{PHCA-DP-SGD (Ours)} & $48.82 \pm 0.95$ & $\mathbf{46.73 \pm 0.42}$ & $\mathbf{2.09}$ \\
\bottomrule
\end{tabular}
\end{table}

PHCA-DP-SGD achieves $46.73\%$ test accuracy after compression, representing a $3.48$ percentage point improvement over standard DP-SGD ($43.25\%$) and a $2.16$ point improvement over AdaDPIGU. Notably, PHCA-DP-SGD matches pre-pruning performance ($46.73\%$ vs. $46.91\%$) \textbf{without requiring public data}.

\subsection{Privacy-Utility Trade-off}

Figure~\ref{fig:privacy_utility} shows the privacy-utility trade-off across different privacy budgets ($\epsilon \in \{1, 3, 8\}$).

\begin{figure}[t]
\centering
\includegraphics[width=0.8\linewidth]{figures/figure1_privacy_utility.pdf}
\caption{Privacy-utility trade-off on CIFAR-10 with 70\% sparsity. PHCA-DP-SGD consistently outperforms standard DP-SGD across all privacy budgets.}
\label{fig:privacy_utility}
\end{figure}

At $\epsilon=1$, PHCA-DP-SGD achieves $38.5\%$ vs. $35.2\%$ for standard DP-SGD. At $\epsilon=8$, the gap narrows but PHCA-DP-SGD still maintains a $2.4$ point advantage ($54.2\%$ vs. $51.8\%$).

\subsection{Compression Robustness}

Figure~\ref{fig:sparsity} examines performance across different compression ratios.

\begin{figure}[t]
\centering
\includegraphics[width=0.8\linewidth]{figures/figure5_sparsity.pdf}
\caption{Test accuracy vs. compression sparsity at $\epsilon=3$. PHCA-DP-SGD shows greater robustness to aggressive compression.}
\label{fig:sparsity}
\end{figure}

PHCA-DP-SGD maintains higher accuracy across all sparsity levels, with the advantage increasing at higher compression ratios. This demonstrates that compression-aware training produces models whose important parameters are more robust to pruning.

\subsection{Ablation Study}

Table~\ref{tab:ablation} presents the ablation study, systematically removing each component of PHCA-DP-SGD.

\begin{table}[t]
\caption{Ablation study on CIFAR-10 at $\epsilon=3$, 70\% sparsity. Each row removes one component from the full method.}
\label{tab:ablation}
\centering
\begin{tabular}{lc}
\toprule
Configuration & Post-Comp. Acc. ($\%$) $\uparrow$ \\
\midrule
Full PHCA-DP-SGD & $\mathbf{46.73 \pm 0.42}$ \\
No per-parameter clipping & $44.85 \pm 0.65$ \\
No compression-aware noise & $45.12 \pm 0.58$ \\
Binary masking (vs. continuous) & $44.57 \pm 0.71$ \\
Fixed survival probabilities & $43.92 \pm 0.82$ \\
\bottomrule
\end{tabular}
\end{table}

Both per-parameter clipping and compression-aware noise contribute significantly to performance. The largest drop ($2.81$ points) occurs when replacing continuous weights with binary masking, validating our design choice.

\subsection{Hyperparameter Sensitivity}

Figure~\ref{fig:hyperparam} shows accuracy as a function of $\alpha$ (discount factor) and $\beta$ (noise reduction rate).

\begin{figure}[t]
\centering
\includegraphics[width=0.8\linewidth]{figures/figure4_hyperparameter.pdf}
\caption{Hyperparameter sensitivity: test accuracy vs. $(\alpha, \beta)$ at $\epsilon=3$, 70\% sparsity. The optimal region is around $\alpha=0.1$, $\beta=0.5$.}
\label{fig:hyperparam}
\end{figure}

The method is relatively robust to hyperparameter choices, with a broad optimal region around $\alpha=0.1$, $\beta=0.5$. Notably, $\alpha=1.0$ (uniform clipping) performs poorly, confirming that compression-aware weighting is essential.

\subsection{Method Comparison Summary}

Figure~\ref{fig:comparison} provides a visual comparison of all methods.

\begin{figure}[t]
\centering
\includegraphics[width=0.8\linewidth]{figures/figure2_method_comparison.pdf}
\caption{Method comparison at $\epsilon=3$, 70\% sparsity. Error bars show standard deviation across 3 random seeds.}
\label{fig:comparison}
\end{figure}

\subsection{Training Dynamics}

Figure~\ref{fig:training} shows training curves for PHCA-DP-SGD vs. standard DP-SGD.

\begin{figure}[t]
\centering
\includegraphics[width=0.8\linewidth]{figures/figure6_training_curves.pdf}
\caption{Training curves: test accuracy over epochs at $\epsilon=3$. PHCA-DP-SGD maintains higher accuracy throughout training.}
\label{fig:training}
\end{figure}

PHCA-DP-SGD converges faster and achieves higher final accuracy, suggesting that compression-aware training also improves optimization dynamics.

\section{Discussion and Limitations}

\subsection{Key Findings}

Our experiments confirm that compression-aware DP training significantly improves privacy-utility trade-offs for models that will undergo post-hoc compression. The key insight---that privacy budget should be allocated based on parameter survival probability---is validated by the $3.48$ percentage point improvement over standard DP-SGD.

Surprisingly, PHCA-DP-SGD matches pre-pruning performance without requiring public data. This suggests that the benefit of pre-pruning comes primarily from concentrating updates on important parameters, which can be achieved through continuous clipping weights rather than binary selection.

\subsection{Limitations}

\textbf{Theoretical guarantees.} While we derive a privacy bound for PHCA-DP-SGD, the tightness of this bound depends on the minimum clipping weight $w_{\min} = \alpha$. Tighter analysis accounting for the distribution of survival probabilities could yield better privacy-utility trade-offs.

\textbf{Compression types.} We focus on magnitude pruning; other compression methods (quantization, structured pruning, knowledge distillation) may require different survival probability estimators.

\textbf{Scale.} Our experiments are on CIFAR-10 with ResNet-18. Scaling to larger datasets and models remains future work.

\textbf{Computational overhead.} PHCA-DP-SGD requires maintaining per-parameter gradient EMAs, adding modest memory overhead ($\sim$20\% for ResNet-18).

\subsection{Broader Impact}

This research addresses a critical gap in deploying privacy-preserving ML systems. By enabling better privacy-utility trade-offs without requiring public data, we expand the scenarios where DP can be practically deployed. However, we caution that improved utility should not be interpreted as permission to deploy in high-stakes domains without thorough privacy auditing.

\section{Conclusion}

We presented PHCA-DP-SGD, the first method to optimize DP training for post-hoc compression without requiring public data. By using continuous per-parameter clipping weights based on estimated survival probabilities, we achieve $3.48$ percentage points higher accuracy than standard DP-SGD after compression on CIFAR-10, matching pre-pruning performance without its public data requirement. Our work demonstrates that compression and privacy should be jointly optimized, not treated as separate concerns.

Future work includes: (1) extending to other compression methods (quantization, distillation); (2) tighter privacy analysis accounting for the full distribution of survival probabilities; (3) scaling to larger models and datasets; (4) federated learning settings where compression is applied at edge devices.

\bibliographystyle{plainnat}
\begin{thebibliography}{10}

\bibitem[Abadi et~al.(2016)]{abadi2016deep}
Martin Abadi, Andy Chu, Ian Goodfellow, H.~Brendan McMahan, Ilya Mironov, Kunal Talwar, and Li~Zhang.
\newblock Deep learning with differential privacy.
\newblock In \emph{Proceedings of the 2016 ACM SIGSAC Conference on Computer and Communications Security}, pages 308--318, 2016.

\bibitem[Adamczewski et~al.(2023)]{adamczewski2023prepruning}
Kamil Adamczewski, Yingchen He, and Mijung Park.
\newblock Pre-pruning and gradient-dropping improve differentially private image classification.
\newblock \emph{arXiv preprint arXiv:2306.11754}, 2023.

\bibitem[Adamczewski \& Park(2023)]{adamczewski2023differential}
Kamil Adamczewski and Mijung Park.
\newblock Differential privacy meets neural network pruning.
\newblock \emph{arXiv preprint arXiv:2303.04612}, 2023.

\bibitem[Balle et~al.(2018)]{balle2018privacy}
Borja Balle, Gilles Barthe, and Marco Gaboardi.
\newblock Privacy amplification by subsampling: Tight analyses via couplings and divergences.
\newblock In \emph{Advances in Neural Information Processing Systems}, 2018.

\bibitem[Dwork et~al.(2006)]{dwork2006calibrating}
Cynthia Dwork, Frank McSherry, Kobbi Nissim, and Adam Smith.
\newblock Calibrating noise to sensitivity in private data analysis.
\newblock In \emph{Theory of Cryptography Conference}, pages 265--284, 2006.

\bibitem[He et~al.(2016)]{he2016deep}
Kaiming He, Xiangyu Zhang, Shaoqing Ren, and Jian Sun.
\newblock Deep residual learning for image recognition.
\newblock In \emph{Proceedings of the IEEE Conference on Computer Vision and Pattern Recognition}, pages 770--778, 2016.

\bibitem[Hu et~al.(2024)]{hu2024past}
Qiang Hu, Yufei Chen, Jasper Lai, Chao Chen, Zhongfeng Niu, Zhi Zhang, Yansong Gao, Anmin Fu, and Jinghua Jiang.
\newblock Defending membership inference attacks via privacy-aware sparsity tuning.
\newblock \emph{arXiv preprint arXiv:2410.06814}, 2024.

\bibitem[Krizhevsky(2009)]{krizhevsky2009learning}
Alex Krizhevsky.
\newblock Learning multiple layers of features from tiny images.
\newblock Technical report, University of Toronto, 2009.

\bibitem[Li et~al.(2025)]{li2025compleak}
Na Li, Yansong Gao, Hongsheng Hu, Boyu Kuang, and Anmin Fu.
\newblock CompLeak: Deep learning model compression exacerbates privacy leakage.
\newblock \emph{arXiv preprint arXiv:2507.16872}, 2025.

\bibitem[Mironov(2017)]{mironov2017renyi}
Ilya Mironov.
\newblock R\'{e}nyi differential privacy.
\newblock In \emph{2017 IEEE 30th Computer Security Foundations Symposium (CSF)}, pages 263--275, 2017.

\bibitem[Shokri et~al.(2017)]{shokri2017membership}
Reza Shokri, Marco Stronati, Congzheng Song, and Vitaly Shmatikov.
\newblock Membership inference attacks against machine learning models.
\newblock In \emph{2017 IEEE Symposium on Security and Privacy (SP)}, pages 3--18. IEEE, 2017.

\bibitem[Yousefpour et~al.(2021)]{yousefpour2021opacus}
Ashkan Yousefpour, Igor Shilov, Alexandre Sablayrolles, Davide Testuggine, Karthik Prasad, Mani Malek, John Nguyen, Sayan Ghosh, Akash Bharadwaj, Jessica Zhao, Graham Cormode, and Ilya Mironov.
\newblock Opacus: User-friendly differential privacy library in {PyTorch}.
\newblock \emph{arXiv preprint arXiv:2109.12298}, 2021.

\bibitem[Zhang \& Xie(2025)]{zhang2025adadpigu}
Huiqi Zhang and Fang Xie.
\newblock AdaDPIGU: Differentially private SGD with adaptive clipping and importance-based gradient updates for deep neural networks.
\newblock \emph{arXiv preprint arXiv:2507.06525}, 2025.

\end{thebibliography}

\end{document}
